
% --------------------------------------------------------------- CONFIGURATIONS

\documentclass[a4paper,12pt,final,oneside]{book}

\usepackage{rapport}


% -------------------------------------------------------------- META: CONSTANTS

\newcommand{\reporttitle}{Informatique D\'ecisionnelle}
\newcommand{\enseignants}{M~\textsc{Miquel}\\A~\textsc{Tchounikine}}
\newcommand{\reportauthor}{Guillaume~\textsc{Abadie}\\Thierry~\textsc{Cantenot}\\Juliette~\textsc{Courlet}\\Rémi~\textsc{Domingues}\\Adrien~\textsc{Duffy-Coissard}\\Ahmed~\textsc{Kachkach}}
\newcommand{\reportsubject}{Livrable de projet}
\newcommand{\stagetopic}{Mortality Data Pack}
\newcommand{\dateperiod}{du 7 avril au 17 avril 2014}
\newcommand{\HRule}{\rule{\linewidth}{0.5mm}}
\setlength{\parskip}{1ex} % Espace entre les paragraphes

\hypersetup{
	pdftitle={\reporttitle},%
		pdfauthor={\reportauthor},%
		pdfsubject={\reportsubject},%
		pdfkeywords={INSA Lyon} {Mortality Data Pack} {Informatique decisionnelle}
}

\title{\reporttitle}
\author{\reportauthor}
%\setcounter{tocdepth}{4}


% ------------------------------------------------------------------------- FILE

\begin{document}


    % ------------------------------------------------------------------- HEADER

	\renewcommand{\chaptername}{} %\renewcommand{\thechapter}{}
	\renewcommand{\contentsname}{Sommaire}

	\pagestyle{empty}
	\pagenumbering{Roman}


    % ------------------------------------------------------------ HEADER: TITLE

	% Inspiré de http://en.wikibooks.org/wiki/LaTeX/Title_Creation
\begin{center}
	\begin{minipage}[t]{0.48\textwidth}
	  \begin{flushleft}
	    \includegraphics [width=40mm]{images/logo_INSA.png} \\[0.5cm]
			INSA Lyon\\
			20, avenue Albert Einstein\\
			69621 Villeurbanne Cedex
	  \end{flushleft}
	\end{minipage}
	\begin{minipage}[t]{0.48\textwidth}
	  \begin{flushright}
	    %\includegraphics [width=60mm]{images/logo_Passau.jpg} \\[0.5cm]
	    %Universität Passau\\
		%Innstraße, 3\\
		%	D-94032 Passau
	  \end{flushright}
	\end{minipage} \\[2cm]

	\textsc{\Large \reportsubject}\\[0.3cm]
	\HRule \\[0.4cm]
	{\Huge \bfseries \reporttitle}\\[0.3cm]
	{\LARGE \bfseries «~\stagetopic~»}\\[0.3cm]
	{\Large \dateperiod}\\[0.4cm]
	\HRule \\[1cm]

	\includegraphics [scale=0.35]{images/application-xml.png} \\[0.7cm]
	\begin{minipage}[t]{0.4\textwidth}
	  \begin{flushleft} \large
	    \emph{Hexanôme~:}\\
	    \small \reportauthor
	  \end{flushleft}
	\end{minipage}
	\begin{minipage}[t]{0.5\textwidth}
	  \begin{flushright} \large
	    \emph{Enseignants~:} \\
	    \enseignants
	  \end{flushright}
	\end{minipage}

	\vfill
	\footnotesize{Année scolaire 2013-2014}
\end{center}



    % --------------------------------------------------- HEADER: CONFIGURATIONS

	\sloppy          % Justification moins stricte : des mots ne dépasseront pas des paragraphes

    \frontmatter
		\pagestyle{empty}
		\tableofcontents
		\addtocontents{toc}{\protect\thispagestyle{empty}}

	\mainmatter
	\pagestyle{headings}

	\renewcommand{\chaptermark}[1]{\markboth{\MakeUppercase{\chaptername\ \thechapter.\ #1}}{}}
	\renewcommand{\sectionmark}[1]{\markright{\thesection{} #1}}


    % ------------------------------------------------------------------ CONTENT

	\chapter{Analyse des données}

\section{Visualisation des données}

Requête utile pour voir la répartition des causes de décés selon le sexe :

\begin{lstlisting}[frame=single, language=SQL]
SELECT sexage.SEX, causes.LABEL, COUNT(*) AS Number
FROM london INNER JOIN
        sexage ON london.SEXAGE = sexage.SEXAGE INNER JOIN
        causes on london.CAUSE = causes.CAUSE
GROUP BY sexage.sex, causes.label
\end{lstlisting}

\section{Séparation des données de sexe et d'age}

    Vu le format des données de la table SexAge, dissocier l'effet du sexe et de l'âge s'avérait hardu.

    Nous avons donc modifié la table en créant deux nouvelles colonnes: Sex, nvarchar qui contiendra le sexe (M/F), et Age,
    nvarchar qui contiendra la tranche d'âge (<1, 10-14, ...).

    Nous avons alors lancé la requête suivante afin de séparer les données de ``sexage'' sur ces deux colonnes :

    \begin{lstlisting}[frame=single, language=SQL]
UPDATE sexage
SET Sex=LEFT(label, 1), Age=RIGHT(label, LEN(label) - 1);
    \end{lstlisting}

\section{Incohérences dans les nombre de morts}

    En analysant la table des morts de chaque région, nous avons remarqué que la somme des morts par cause de niveau 2 était
    inférieure au nombre de morts par cause de niveau 1 (càd de toutes les causes de mort).

    Cette incohérence peut être montrée sur la table ``westmids'' (par exemple) avec la requête suivante :

    \begin{lstlisting}[frame=single, language=SQL]
SELECT causes.NIVEAU, SUM(westmids.DEATHS) AS TotalDeaths
FROM westmids INNER JOIN causes ON westmids;CAUSE = causes.CAUSE
GROUP BY causes.NIVEAU
HAVING (causes.NIVEAU = 1) OR (causes.NIVEAU = 2)
    \end{lstlisting}

    Ceci est probablement dû à une erreur lors du remplissage des formulaires (cause de niveau 2 non-remplie).

    La solution qu'on propose est de rajouter une nouvelle cause de mort de niveau 2 représentant une cause indéfinie, afin
    de ne pas sous-estimer le nombre de morts en ignorant ceux dont la cause de niveau 2 n'est pas renseignée.

    \begin{lstlisting}[frame=single, language=SQL]
METTRE LES REQUETES QUI VONT BIEN ICI
    \end{lstlisting}

\section{Union des données dans une nouvelle table}

    TODO : Mettre du blabla aussi.

    \begin{lstlisting}[frame=single, language=SQL]
INSERT INTO deaths
SELECT *
FROM
(
    SELECT * FROM london UNION
    SELECT * FROM mersey UNION
    SELECT * FROM yorkhumb UNION
    SELECT * FROM ntheast UNION
    SELECT * FROM nthwest UNION
    SELECT * FROM stheast UNION
    SELECT * FROM sthwest UNION
    SELECT * FROM eastern UNION
    SELECT * FROM eastminds UNION
    SELECT * FROM wales UNION
    SELECT * FROM westmids UNION
)
    \end{lstlisting}


\section{Visualisation des données}

    Nous avons exploré les données mises à notre disposition, en essayant de trouver des motifs, incohérences et autres motifs particuliers
    dans ces dernières.

    Pour cela, nous avons effectué plusieurs requêtes SQL, comme cette requête qui donne la répartition des décès par sexe pour chaque
    cause.

    \begin{lstlisting}[frame=single, language=SQL]
SELECT sexage.SEX, causes.LABEL, COUNT(*) AS Number
FROM london INNER JOIN
        sexage ON london.SEXAGE = sexage.SEXAGE INNER JOIN
        causes on london.CAUSE = causes.CAUSE
GROUP BY sexage.sex, causes.label
    \end{lstlisting}

    On a donc découvert que, surprise !, il n'y a pas de décès ayant pour cause un avortement chez le sexe masculin.

	%\renewcommand{\chaptermark}[1]{\markboth{\MakeUppercase{#1}}{}}
	%\renewcommand{\sectionmark}[1]{\markright{#1}}

	%\addcontentsline{toc}{part}{Annexes}
	%\part*{Annexes}
	%\appendix
	%\include{implementationExercices}


    % ------------------------------------------------------------------- FOOTER
\end{document}
